\chapter{Repository Structure}
\label{ch:repo}

\section{The tree}

The repository has 118 tracked files (excluding the lockfile). Figure~\ref{fig:tree} shows every meaningful path, with line counts for source files where they help you judge where the weight is.

\begin{figure}[htbp]
\begin{lstlisting}[language={},numbers=none,basicstyle=\ttfamily\scriptsize,frame=single]
VidyaRAG/
|-- .github/workflows/        ci.yml (85) docker.yml (111) eval.yml (77)
|-- app/                      app.py (208) requirements.txt README_space.md
|-- config/
|   |-- default.yaml          base pipeline configuration (95)
|   `-- profiles/             baseline rerank decompose corrective guarded
|-- docs/                     DESIGN.md (395) EVALUATION.md (798) architecture.mmd
|-- eval/
|   |-- goldset/              goldset_v1.jsonl (58 questions) REVIEW.md
|   `-- results/              5 run files: <profile>__<run_id>.json + .md
|-- scripts/                  deploy_space.py (171)
|-- src/vidyarag/
|   |-- __init__.py           __version__ = "1.2.1"
|   |-- _compat.py            StrEnum / UTC shims for Python 3.10
|   |-- settings.py (267)     Settings, PipelineConfig, profile loader
|   |-- cli.py (818)          Typer command-line tool
|   |-- pipeline.py (287)     Pipeline orchestrator
|   |-- api/                  main.py (70) models.py (161) routes_v1.py (194)
|   |-- correct/              grader.py (233) policy.py (104) loop.py (172)
|   |-- evaluation/           goldset (194) retrieval (131) abstention (158)
|   |                         answer_cache (120) runner (610) metrics (425)
|   |                         report (308) verify (596) draft (479)
|   |-- generate/             prompts.py (81) answer.py (128) citations.py (148)
|   |-- guard/                patterns.py (148) input_guard.py (58) context_guard.py (79)
|   |-- ingest/               corpus (128) download (238) parse (377) chunk (290) pipeline (217)
|   |-- llm/                  provider.py (144)
|   |-- observe/              trace.py (199) logging.py (60)
|   |-- retrieve/             dense.py (124) rerank.py (97) decompose.py (169)
|   `-- store/                client.py (69) collection.py (190)
|-- tests/                    28 test modules + conftest.py (4,066 lines)
|-- .dockerignore .env.example .gitattributes .gitignore .pre-commit-config.yaml
|-- ATTRIBUTION.md  LICENSE  Makefile  README.md
|-- Dockerfile  docker-compose.yml  pyproject.toml  uv.lock
`-- (not tracked) data/raw, data/index, .env, .eval_cache/, .venv/
\end{lstlisting}
\caption{Repository layout at commit \code{fa5164b}. Source: 8,549 lines; tests: 4,066 lines.}
\label{fig:tree}
\end{figure}

\paragraph{What is \emph{not} in the repository, on purpose.} The raw PDFs (\code{data/raw}), the built index (\code{data/index}), the \code{.env} secrets file, and all caches are listed in \code{.gitignore}. The comments there explain why: the PDFs are large and reproducible from a checksummed manifest; the index ships to the Space through the Space's own repository; and a token pasted into a scratch file is ``one \code{git add -A} away from history'' \ev{.gitignore}. Personal notes matching \code{NOTES-*.md} are also ignored.

\section{Directory by directory}

\subsection{\texorpdfstring{\code{src/vidyarag/}}{src/vidyarag/} --- the package}
All runtime and evaluation logic. It uses the ``src layout'': the package is not importable from the repository root unless installed (which \code{uv sync} does) or put on \code{sys.path} (which the Space bootstrap does). The evaluation code lives in \code{vidyarag.evaluation} rather than a top-level \code{eval} package because a package named \code{eval} would shadow Python's built-in \code{eval} \ev{src/vidyarag/evaluation/__init__.py:1-6}.

\subsection{\texorpdfstring{\code{config/}}{config/} --- behaviour as data}
\code{default.yaml} holds every tunable value with extensive comments explaining each choice; the five files in \code{profiles/} overlay small changes. Chapter~\ref{ch:config} covers the merge and validation rules.

\subsection{\texorpdfstring{\code{app/}}{app/} --- the Hugging Face Space}
\code{app.py} is the Gradio UI; \code{requirements.txt} is what the Space installs (it deliberately omits Gradio, see Chapter~\ref{ch:deploy}); \code{README_space.md} becomes the Space's \code{README.md}, whose YAML front matter pins the Gradio SDK version.

\subsection{\texorpdfstring{\code{eval/}}{eval/} --- evidence, not code}
The gold set and the committed run results. The comment in \code{.gitignore} states that \code{eval/results/*.json} is committed on purpose: ``the run history is the evidence behind the README results table''; only \code{eval/results/scratch/} is ignored.

\subsection{\texorpdfstring{\code{tests/}}{tests/} --- 28 modules}
Hermetic by default: in-memory Qdrant, fake keys, mocked HTTP. Three live tests are marked \code{costly}. Chapter~\ref{ch:testing} walks through them.

\subsection{\texorpdfstring{\code{docs/}}{docs/} --- design and evaluation narratives}
\code{DESIGN.md} is a decision log with rejected alternatives; \code{EVALUATION.md} is the methodology and results narrative, including a ``What did not work'' section. \code{architecture.mmd} is the Mermaid diagram also embedded in the README; a test fails if the two copies drift \ev{tests/test_docs_consistency.py}.

\subsection{\texorpdfstring{\code{scripts/}}{scripts/} and \texorpdfstring{\code{.github/}}{.github/}}
One deployment script, and three GitHub Actions workflows (CI, Docker, manual evaluation).

\section{File by file}

The tables below cover every important file. ``Runs when'' says what triggers the code; ``If removed/changed'' describes the consequence so you understand why the file matters.

\subsection{Package root and configuration}

\begin{longtable}{@{}>{\raggedright\arraybackslash}p{3.0cm}>{\raggedright\arraybackslash}p{12.5cm}@{}}
\toprule
\textbf{File} & \textbf{Role, key symbols, dependencies, when it runs, why it matters}\\
\midrule
\endhead
\code{__init__.py} & Defines \code{__version__ = "1.2.1"}, the \emph{only} place the version is written. \code{pyproject.toml} reads it via Hatch's dynamic version and the API reports it in OpenAPI. It is a literal rather than \code{importlib.metadata} because the Space imports the package from \code{src/} without installing it. \textbf{If changed:} the CLI \code{version} command, \code{/docs} and the installed metadata all change together; a Docker CI step checks the CLI and API agree.\\
\code{_compat.py} & Exports \code{StrEnum} and \code{UTC}. On Python 3.11+ these are the standard-library objects; on 3.10 a \code{str}+\code{Enum} backport with \code{__str__ = str.__str__}. Imported by \code{settings.py}, \code{correct/}, \code{guard/patterns.py}, \code{evaluation/goldset.py}, and modules using timestamps. \textbf{If removed:} the Hugging Face Space (Python 3.10) crashes on import --- which is exactly what happened before it existed.\\
\code{settings.py} & Two configuration layers. \code{Settings} (pydantic-settings) reads environment variables and \code{.env}: Gemini key, Qdrant mode/URL/key/path/collection, profile name (default \code{guarded}), logging. \code{PipelineConfig} and its nested \code{ChunkingConfig}, \code{RetrievalConfig}, \code{CorrectiveConfig}, \code{GuardrailConfig} describe behaviour, all with \code{extra="forbid"}. \code{load_pipeline_config(profile)} deep-merges \code{default.yaml} with a profile. Also \code{resolve_config_dir()} (env override \code{VIDYARAG_CONFIG_DIR}) and \code{REPO_ROOT}. Imported by almost everything. \textbf{Runs} at every startup. \textbf{If changed carelessly:} a typo in a profile key is a startup error, by design.\\
\code{config/default.yaml} & Every pipeline default, each justified in comments: pinned Gemini models (and why lite), local BGE embeddings, chunk 512/64, retrieve 20 $\rightarrow$ context 5, reranker model, reserved \code{sparse_model}, corrective thresholds 0.8/0.5/2, guardrails off. \textbf{If changed:} every profile inherits the change, which would silently alter the frozen baseline --- a test asserts the baseline keeps all enhancements off.\\
\code{config/profiles/*.yaml} & \code{baseline} (frozen control), \code{rerank}, \code{decompose}, \code{corrective}, \code{guarded} (shipped). Each file's comments record what the ablation was expected to move. Read by \code{load_pipeline_config}.\\
\bottomrule
\end{longtable}

\subsection{Ingestion (\texorpdfstring{\code{src/vidyarag/ingest/}}{ingest})}

\begin{longtable}{@{}>{\raggedright\arraybackslash}p{3.0cm}>{\raggedright\arraybackslash}p{12.5cm}@{}}
\toprule
\textbf{File} & \textbf{Role, key symbols, dependencies, when it runs, why it matters}\\
\midrule
\endhead
\code{corpus.py} & \code{BookSpec} (frozen dataclass: slug, title, edition, dates, URLs, UUID, ISBN, licence) and the two instances \code{BIOLOGY} and \code{ANATOMY_AND_PHYSIOLOGY}; \code{CORPUS} tuple; \code{get_book(slug)}. Explains why only these first editions are CC~BY (second editions are CC~BY-NC-SA). Used by download, parse, ingest pipeline and CLI. \textbf{If changed:} adding a book changes chunk ids and every evaluation number; the licence must be re-verified per edition.\\
\code{download.py} & \code{download_book}, \code{download_corpus}, \code{load_manifest}; \code{DownloadRecord}, \code{Manifest}. Streams with \code{httpx}, resumes with HTTP \code{Range}, writes to \code{.part} then renames, handles 416 and ignored-Range 200 responses, hashes with SHA-256, opens the PDF with PyMuPDF to count pages. Runs on \code{vidyarag download}. \textbf{Why it matters:} the manifest pins \emph{which bytes} produced the index.\\
\code{parse.py} & \code{load_outline}, \code{build_structure_map}, \code{extract_pages}, \code{printed_page_number}, \code{clean_text}, \code{strip_end_matter}, \code{is_question_page}; dataclasses \code{OutlineEntry}, \code{PageText}. Uses PyMuPDF. Runs during ingestion. \textbf{Why it matters:} every citation's chapter, section and printed page come from here; a mistake here makes citations wrong while everything downstream looks fine.\\
\code{chunk.py} & \code{Chunk}, \code{split_sentences}, \code{chunk_pages} and the private \code{_pack}/\code{_carry_overlap}. Uses \code{count_tokens} from the provider (BGE tokeniser). Runs during ingestion. \textbf{Why it matters:} chunk boundaries decide what can be retrieved; the ``what is tissue?'' weakness is a chunk-boundary effect.\\
\code{pipeline.py} & \code{ingest(client, ...)}, \code{chunks_for_book}, \code{IngestReport}, \code{BookIngestStats}. Orchestrates per book: parse, chunk, skip existing ids, embed in batches, upsert, write report. Called by \code{vidyarag ingest}. \textbf{Note:} the report path defaults to the fixed \code{data/ingest_run.json} regardless of which index was built (Section~\ref{sec:finding-ingest-report}).\\
\bottomrule
\end{longtable}

\subsection{Store, models and observability}

\begin{longtable}{@{}>{\raggedright\arraybackslash}p{3.0cm}>{\raggedright\arraybackslash}p{12.5cm}@{}}
\toprule
\textbf{File} & \textbf{Role, key symbols, dependencies, when it runs, why it matters}\\
\midrule
\endhead
\code{store/client.py} & \code{build_client(settings)} returns a \code{QdrantClient} for embedded (in-memory or folder), server or cloud mode; \code{describe_target} prints the target with the API key redacted. \textbf{Why it matters:} the only place that knows how Qdrant is reached, so tests, demo and evaluation differ only by environment variables.\\
\code{store/collection.py} & \code{POINT_NAMESPACE}, \code{point_id} (uuid5), \code{DENSE_VECTOR = "dense"}, \code{ensure_collection} (cosine, dimension check, \code{book_slug} payload index), \code{build_payload}, \code{make_points}, \code{upsert_points}, \code{count_points}, \code{existing_chunk_ids}. Used by ingestion, retrieval and the API health route.\\
\code{llm/provider.py} & \code{get_embedder} and \code{get_tokenizer} (cached), \code{count_tokens}, \code{embed_texts}, \code{get_gemini_client} (90-second timeout, lazy SDK import, clear error if no key), \code{REQUEST_TIMEOUT_MS}. \textbf{Why it matters:} the swap point for every model.\\
\code{observe/trace.py} & \code{QueryTrace} (stages, usage, chunk-id lists, guard and loop records), \code{StageTiming}, \code{Usage}, \code{price_call}, \code{GEMINI_PRICING}, \code{Timer}. Every stage writes into it. \textbf{Known issue:} totals double-count nested stages (Section~\ref{sec:latency-bug}).\\
\code{observe/logging.py} & \code{configure_logging}, \code{get_logger} (structlog, JSON or console). Called only by the API's startup.\\
\bottomrule
\end{longtable}

\subsection{Query stages}

\begin{longtable}{@{}>{\raggedright\arraybackslash}p{3.0cm}>{\raggedright\arraybackslash}p{12.5cm}@{}}
\toprule
\textbf{File} & \textbf{Role, key symbols, dependencies, when it runs, why it matters}\\
\midrule
\endhead
\code{retrieve/dense.py} & \code{RetrievedChunk} (frozen dataclass with all citation fields and \code{prior_score}), \code{retrieve_dense} (embed query, optional \code{book_slug} filter, \code{query_points} on the \code{dense} vector). The frozen baseline retriever; also used by the \code{/v1/search} endpoint, reranker, decomposer and gold-set tooling.\\
\code{retrieve/rerank.py} & \code{get_reranker} (cached fastembed \code{TextCrossEncoder}), \code{rerank}, \code{rank_movement}. Runs when \code{use_reranker} is true.\\
\code{retrieve/decompose.py} & \code{Decomposition} schema, \code{DECOMPOSE_PROMPT}, \code{decompose} (falls back to ``not multi-hop'' on any failure), \code{reciprocal_rank_fusion} ($k=60$), \code{retrieve_decomposed}. Runs only in the \code{decompose} profile; measured and rejected but kept so the negative result stays reproducible.\\
\code{generate/prompts.py} & \code{Prompt}, \code{ANSWER_V1_SYSTEM}, \code{ANSWER_V1_USER}, \code{ANSWER_PROMPT_VERSION = "answer-v1"}, \code{build_answer_prompt}, \code{NO_CONTEXT_ANSWER}. \textbf{If the prompt text changes}, the version must be bumped, because the version is part of the evaluation cache key and the trace.\\
\code{generate/citations.py} & \code{Citation}, \code{format_context}, \code{extract_markers}, \code{resolve_citations}, \code{strip_invalid_markers}, \code{render_references}; the \code{_MARKER_GROUP} regex. The component that prevents fabricated citations from reaching a reader.\\
\code{generate/answer.py} & \code{GeneratedAnswer}, \code{generate_answer}, \code{_extract_text}, \code{_record_usage}. Calls Gemini with a system instruction; never calls it with empty context.\\
\code{correct/grader.py} & \code{ClaimVerdict}, \code{Claim}, \code{GradedAnswer} (JSON schema), \code{Groundedness} (score, \code{measured}, \code{failed_claims}, \code{refuses}), \code{GRADER_PROMPT}, \code{grade_answer}.\\
\code{correct/policy.py} & \code{Decision}, \code{CorrectivePolicy.decide}; validates threshold ordering. The most consequential numbers in the system live here.\\
\code{correct/loop.py} & \code{ABSTENTION_TEXT}, \code{Attempt}, \code{LoopOutcome}, \code{reformulate}, \code{run_corrective_loop}.\\
\code{guard/patterns.py} & \code{Category}, \code{Detection}, the compiled regex families and \code{scan}. The false-positive engineering lives in its comments.\\
\code{guard/input_guard.py}, \code{guard/context_guard.py} & \code{screen_input} with \code{REFUSAL}; \code{screen_context} with \code{ContextVerdict}.\\
\code{pipeline.py} & \code{Answer}, \code{Pipeline} (\code{retrieve}, \code{answer}, \code{_answer_corrective}, \code{close}, lazy \code{llm}), \code{build_pipeline}. The heart of query time.\\
\bottomrule
\end{longtable}

\subsection{Interfaces, evaluation and deployment}

\begin{longtable}{@{}>{\raggedright\arraybackslash}p{3.0cm}>{\raggedright\arraybackslash}p{12.5cm}@{}}
\toprule
\textbf{File} & \textbf{Role, key symbols, dependencies, when it runs, why it matters}\\
\midrule
\endhead
\code{api/main.py} & \code{lifespan} (build pipeline once, log startup, close on shutdown), \code{create_app}, module-level \code{app}. Run by \code{uvicorn vidyarag.api.main:app}.\\
\code{api/models.py} & Pydantic request/response schemas: \code{QueryRequest}, \code{QueryResponse}, \code{CitationOut}, \code{RetrievedOut}, \code{TraceOut}, \code{StageOut}, \code{HealthResponse}, \code{ConfigResponse}, \code{SearchRequest}, \code{SearchHit}, \code{SearchResponse}.\\
\code{api/routes_v1.py} & \code{POST /v1/query}, \code{POST /v1/search}, \code{GET /v1/health}, \code{GET /v1/config}; \code{get_pipeline} dependency.\\
\code{cli.py} & Typer app: \code{version}, \code{config}, \code{download}, \code{ingest}, \code{ask}, \code{health}, \code{eval}, \code{report}, and \code{goldset draft|check|unanswerable|triage}. Forces UTF-8 output on Windows. 0\% covered by tests.\\
\code{app/app.py} & Gradio \code{build_ui}, \code{ask}, \code{get_pipeline} singleton, renderers, ZeroGPU declaration.\\
\code{evaluation/goldset.py} & \code{QuestionType}, \code{Provenance}, \code{GoldQuestion} (with validators), \code{load_goldset}, \code{summarise_goldset}.\\
\code{evaluation/retrieval.py} & \code{recall_at_k}, \code{hit_at_k}, \code{reciprocal_rank}, \code{RetrievalScores}, \code{score_retrieval}.\\
\code{evaluation/abstention.py} & \code{AbstentionStats}, \code{is_structural_abstention}, \code{judge_abstention}, \code{summarise_abstention}, \code{ABSTENTION_JUDGE_PROMPT}. \textbf{Contains the truncation bug} (Section~\ref{sec:judge-bug}).\\
\code{evaluation/answer_cache.py} & \code{answer_key}, \code{abstention_key}, \code{AnswerCache}.\\
\code{evaluation/metrics.py} & The only RAGAS importer: \code{MetricSuite}, \code{RateLimiter}, \code{SampleScores}, \code{LocalEmbedder}, \code{check_grading_dependencies}, \code{GEMINI_OPENAI_BASE_URL}.\\
\code{evaluation/runner.py} & \code{run_evaluation}, \code{SampleResult}, \code{EvalRun}, \code{latest_run}, \code{profiles_with_runs}, constants for pacing, retries and validity.\\
\code{evaluation/report.py} & \code{render_report}, \code{render_comparison}.\\
\code{evaluation/verify.py} & \code{propose_unanswerable}, \code{check_unanswerable}, \code{in_domain_threshold}, \code{triage_answerable} and prompts; 43\% covered.\\
\code{evaluation/draft.py} & \code{scroll_chunks}, \code{sample_chunks}, \code{find_related_chunk}, \code{draft_factual}, \code{draft_multihop}, review-sheet writer; 0\% covered.\\
\code{scripts/deploy_space.py} & \code{PAYLOAD}, \code{bootstrapped_app}, \code{stage}, \code{main}.\\
\code{Dockerfile}, \code{docker-compose.yml} & Two-stage API image (non-root, HTTP healthcheck); local Qdrant server for development.\\
\code{pyproject.toml} & Metadata, dependency groups (runtime, \code{dev}, \code{eval}), console script, and tool configuration (ruff, black, mypy strict, pytest markers, coverage floor 60\%).\\
\code{Makefile} & Thin \code{uv run} wrappers: \code{install}, \code{health}, \code{download}, \code{ingest}, \code{ui}, \code{serve}, \code{test}, \code{test-cov}, \code{lint}, \code{format}, \code{typecheck}, \code{check}, \code{qdrant-up/down}, \code{clean}.\\
\code{.pre-commit-config.yaml} & gitleaks first, then hygiene hooks, ruff, black, mypy.\\
\bottomrule
\end{longtable}
